\section{Conclusion}
\label{sec:conclusion}

Federal mitigation grants are meant to turn the aftermath of a disaster into an opportunity to reduce future risk, but in HMGP that opportunity becomes a project only when a local government applies. Across two decades of tornado and severe-storm declarations, the smallest third of counties obtained HMGP projects in roughly one fewer case out of every ten county--declarations than larger counties designated in the same declaration, and they were about half as likely to obtain the safe rooms and wind retrofits that address tornado risk directly. The gap opens before any request is reviewed.

For regional economics, the main lesson concerns the delivery of intergovernmental transfers. When federal funds for local public investment must be claimed through applications, the distribution of investment across places reflects how local government is organized and the capacity of individual governments to apply, as well as hazard and need. Part of the small-county gap is associated with the smaller number of local entities in small counties; the rest appears for county governments themselves and in their weaker response to tornado hits. Reimbursing management costs after an award did not detectably change this pattern, while state preparation of applications is associated with a smaller gap. These results are descriptive, and identifying the causal effect of local capacity on grant participation, for example through variation in technical assistance across states or over time, remains a task for future research. The evidence nonetheless indicates that reducing geographic disparities in mitigation investment requires attention to how applications are prepared, and not only to how funds are allocated.
